\bigskip
\bigskip
\subsubsection*{Խնդիրներ և հարցեր թեմա 12-ի վերաբերյալ}

\begin{enumerate}[label=\thesection.\arabic*., series=probset12]

% 12.1
\item Ապացուցեք․ $(\R, \text{սովոր.})$ թվային ուղղի $(-\infty, 0] \cup (0, +\infty)$ տրոհումից ստաց\-վող ֆակտոր-տարածությունը հոմեոմորֆ է կապակցված կետազույգ տարա\-ծու\-թյա\-նը (սահմանումը տե՛ս թեմա 4-ում):

% 12.2
\item Դիտարկենք $(\R, \text{սովոր.})$ տարածության կետերի միջև $R$ երկտեղ հարա\-բե\-րու\-թյուն, ըստ որի $x_1 R x_2 \, \Leftrightarrow \, (x_1 - x_2)$-ը ամբողջ թիվ է: Ապացուցեք․
\begin{enumerate}
    \item[ա)] $R$-ը համարժեքության հարաբերություն է,
    \item[բ)] $\faktor{\R}{R}$ ֆակտոր-տարածությունը հոմեոմորֆ է $(\R, \text{սովոր.})$ տարածության \linebreak շրջա\-նագծի:
\end{enumerate}

\begin{hint}
Օգտվեք թեմա 2-ում \hyperref[ch2ex9]{օրինակ 9}-ից և այս թեմայի \hyperref[ch12ex2]{օրինակ 2}-ից։
\end{hint}

% 12.3
\item Դիտարկենք $\R^2$ կոորդինատային հարթությունը սովորական տոպոլոգիայով և $R$ երկտեղ հարաբերություն նրա կետերի միջև, ըստ որի $(x_1, y_1) R (x_2, y_2) \Leftrightarrow x_1 - x_2$ և $y_1 - y_2$ թվերը ամբողջ են։ Ապացուցեք․
\begin{enumerate}
    \item[ա)] $R$-ը համարժեքության հարաբերություն է,
    \item[բ)] $\faktor{\R^2}{R}$ ֆակտոր-տարածությունը հոմեոմորֆ է տորի:
\end{enumerate}

\begin{hint}
Օգտվեք թեմա 2-ում \hyperref[ch2ex10]{օրինակ 10}-ից:
\end{hint}

% 12.4
\item Ճի՞շտ է արդյոք պնդումը․ եթե ունենք տոպոլոգիական տարածությունների \linebreak ${f: X \to Y}$ անընդհատ, սյուրյեկտիվ արտապատկերում, ապա $Y$-ը $X$-ի ֆակ\-տոր-տարածությունն է որոշված $f$ արտապատկերումով։

\begin{hint}
Դիտարկեք, օրինակ, $(X, \text{դիսկր․})$, $(Y, \text{անտիդ․})$ տա\-րա\-ծու\-թյուն\-նե\-րը, որտեղ $X=\{x_1, x_2\}$, $Y=\{y_1, y_2\}$, և ${f: X \to Y}$, ${f(x_1)=y_1}$, ${f(x_2)=y_2}$ արտապատկերումը։ Ցույց տվեք, որ $Y$ բազմության տվյալ տոպոլոգիան չի համընկնում $f$ արտապատկերումով որոշվող $Y$-ի ֆակտոր-տոպոլոգիայի հետ։
\end{hint}

% 12.5
\item Ապացուցեք, որ $\mathbb{RP}^1$ պրոյեկտիվ ուղիղը հոմեոմորֆ է շրջանագծի:

\begin{hint}
Ներկայացրեք $\mathbb{RP}^1$-ը որպես փակ կիսաշրջանագծի ֆակտոր-տա\-րա\-ծություն:
\end{hint}

% 12.6
\item Ստուգեք․ արդյո՞ք բաց (փակ) արտապատկերում է \hyperref[ch12ex1]{օրինակ 1}-ում ${P: X \to \faktor{X}{R}}$ կանոնական պրոյեկցիան։

% 12.7
\item Ճի՞շտ է արդյոք պնդումը․ ցանկացած $\faktor{X}{R}$ ֆակտոր-տարածության դեպքում $P: X \to \faktor{X}{R}$ կանոնական պրոյեկցիան բաց արտապատկերում է:

\begin{hint}
Դիտարկեք \hyperref[ch12ex2]{օրինակ 2}-ում $[0, \frac{1}{2}) \subset [0, 1]$ ենթաբազմության կեր\-պարը կանոնական պրոյեկցիայի դեպքում:
\end{hint}

% 12.8
\item Ապացուցեք․ կամայական $\faktor{X}{R}$ ֆակտոր-տարածության դեպքում ${P: X \to \faktor{X}{R}}$ կանոնական պրոյեկցիան բաց (փակ) արտապատկերում է այն և միայն այն դեպքում, երբ $X$-ի ցանկացած $U$ բաց ենթաբազմության դեպքում $P^{-1}\big(P(U)\big)$ ենթաբազմությունը բաց է (փակ է) $X$-ում։
\end{enumerate}

\par Մինչև հաջորդ խնդիրներին անցնելը տանք մի սահմանում:

\par Դիցուք $X$-ը տոպոլոգիական տարածություն է, իսկ $A$-ն $X$-ի որևէ ենթաբազ\-մու\-թյուն է։ Դիտարկենք $R$ համարժեքության հարաբերություն $X$-ի վրա համարելով՝
\begin{enumerate}
    \item[ա)] տեղի ունի $a_1 R a_2$ հարաբերություն $A$-ի ցանկացած երկու $a_1, a_2$ տարրերի դեպ\-քում,
    \item[բ)] եթե $x_1, x_2 \in X \setminus A$, կամ $x_1 \in A$ և $x_2 \in X \setminus A$, ապա տեղի ունի $x_1 R x_2$ այն \linebreak և միայն այն դեպքում, երբ $x_1 = x_2$։
\end{enumerate}
\par Այսպիսով համարժեքության դասերը $X$-ի միկետանոց ենթաբազմություններն են և $A$-ն:

\par Այս դեպքում $\faktor{X}{R}$ ֆակտոր-տարածության մասին ասում են, որ այն ստացվում է $X$-ից՝ \textbf{սեղմելով մի կետի} $A$ ենթաբազմությունը։ Նշանակենք այդ տարածու\-թյու\-նը $\faktor{X}{A}$ սիմվոլով:

\begin{enumerate}[label=\thesection.\arabic*., resume=probset12]

% 12.9
\item Դիցուք $X$-ը $\R^2$ կոորդինատային հարթության $\mathcal{B}^2 = \{(x, y) \mid x^2+y^2 \le 1\}$ միավոր շրջանն է, իսկ $A$-ն նրա եզրային $\mathcal{S}^1 = \{(x, y) \mid x^2+y^2 = 1\}$ միավոր շրջանագիծն է։ Ապացուցեք, որ $\faktor{X}{A}$ ֆակտոր-տարածությունը հոմեոմորֆ է $\mathcal{S}^2 = \{(x, y, z) \mid x^2+y^2+z^2=1\}$ սֆերային։
Պատկերավոր ասած՝ շրջանի եզրը կետի սեղմելիս ստացվում է սֆերա։

\begin{hint}
Նախ կառուցեք հոմեոմորֆիզմ $\mathcal{B}^2 = \{(x, y) \mid x^2+y^2 < 1\}$ ան\-եզր շրջանի և առանց մի կետ $\mathcal{S}^2 \setminus \{(0, 0, 1)\}$ շրջանի միջև:
\end{hint}

% 12.10
\item Դիցուք $X$-ը $x^2+y^2=1$ ուղիղ շրջանային գլանն է $\R^3$-ում, իսկ $A$-ն $x^2+y^2=1$, $z=0$ շրջանագիծն է։ Ապացուցեք, որ $\faktor{X}{A}$ ֆակտոր-տարածությունը հոմեոմորֆ է $Y = \{(x, y, z) \mid x^2+y^2-z^2=0\}$ կոնական մակերևույթին։ Պատ\-կե\-րա\-վոր ասած՝ գլանում մի շրջանագիծ կետի սեղմելիս ստացվում է կոնական մակ\-երևույթ:

\begin{hint}
Նախ ցույց տվեք, որ $(x, y, z) \mapsto (xz, yz, z)$ համադրումը որոշում է հոմեոմորֆիզմ $X \setminus A$ և $Y \setminus \{(0, 0, 0)\}$ տարածությունների միջև:
\end{hint}

\end{enumerate}
